%% =============================================================================
%% PAPER 1 - Section 5: Factor Selection and Spanning Regressions
%% =============================================================================
%% TEXT ONLY — no inline tables/figures.
%% Tables/figures are collected in paper1_tables_final.tex and
%% paper1_figures_final.tex (at end of document, per EDHEC format guide).
%% =============================================================================


% ─────────────────────────────────────────────────────────────────────────────
\subsection{Principal Component Analysis}
\label{sec:pca}
% ─────────────────────────────────────────────────────────────────────────────

\subsubsection{Rolling PCA Methodology}
\label{sec:pca_method}

% TODO: Write full prose.
% Z-score standardization of all 71 candidate factors.
% Rolling window W=24 months (Rebonato; Ludvigson \& Ng 2009).
% K=8 fixed PCs (cumulative variance $\approx$ 80\%, Figure~\ref{fig:pca_scree}).
% Predictive (PC_t -> R_{t+1}) and contemporaneous (PC_t -> R_t).
% Spanning regression: R_t = alpha + beta' PC_t + epsilon, OLS + NW HAC.
% Robustness across (W,K) in Appendix~\ref{sec:app_pca_robustness}.


\subsubsection{Results}
\label{sec:pca_results}

% TODO: Write full prose.
% Table~\ref{tab:pca_spanning}: alpha, R-squared adj for 3 strategies.
% Alpha significant for [strategy names]. R-squared moderate.
% Compare with benchmark models (Section~\ref{sec:benchmarks}).


% ─────────────────────────────────────────────────────────────────────────────
\subsection{Adaptive Elastic Net}
\label{sec:aen}
% ─────────────────────────────────────────────────────────────────────────────

\subsubsection{Candidate Factor Universe and Preprocessing}
\label{sec:aen_preproc}

% TODO: Write full prose.
% 71 candidate factors, 5 categories (Appendix~\ref{sec:app_factor_list}).
% Stationarity (ADF), z-score, deduplication $|\rho|>0.95$.


\subsubsection{Estimation and Tuning}
\label{sec:aen_estimation}

% TODO: Write full prose.
% Two-stage AEN (Zou \& Zhang 2009). Stage 1: EN + HQC.
% Adaptive weights $w_j = (|\hat\beta_j^{EN}| + 1/T)^{-\gamma}$, $\gamma=1$.
% Stage 2: weighted $\ell_1$ + uniform $\ell_2$, coordinate descent.


\subsubsection{Stability Selection}
\label{sec:aen_stability}

% TODO: Write full prose.
% Block bootstrap B=100. $\hat\Pi(j, \lambda^*) \geq 0.60$.
% Table~\ref{tab:stability_freq} and Figure~\ref{fig:stability_barplot}.
% Cluster-level aggregation (Faletto \& Bien 2022).


\subsubsection{Post-Selection OLS Results}
\label{sec:aen_results}

% TODO: Write full prose.
% Table~\ref{tab:aen_stable_ols}: alpha, betas, t-stats, R-squared adj.
% Economic interpretation of selected factors.


% ─────────────────────────────────────────────────────────────────────────────
\subsection{Conditional Alpha}
\label{sec:conditional_alpha}
% ─────────────────────────────────────────────────────────────────────────────

% TODO: Write full prose.
% Central result: alpha increases during stress. Limits to arbitrage.
% Table~\ref{tab:conditional_alpha_fs}: Ferson-Schadt with AEN factors.
% Table~\ref{tab:conditional_alpha_threshold}: threshold robustness (AEN).
% Table~\ref{tab:pca_conditional_alpha_fs}: Ferson-Schadt with PCA.
% Table~\ref{tab:pca_conditional_alpha_threshold}: threshold robustness (PCA).
% Robust across factor selection method → structural, not artifact.


% ─────────────────────────────────────────────────────────────────────────────
\subsection{Discussion, Method Comparison, and Robustness}
\label{sec:factor_discussion}
% ─────────────────────────────────────────────────────────────────────────────

% TODO: Write full prose.
% Table~\ref{tab:aen_method_comparison}: AEN vs Adaptive LASSO vs LASSO vs Ridge.
% Alpha robust across all methods. AEN preferred (oracle property).
% Sensitivity and half-sample in Appendix.